\documentclass{article}
\usepackage{amsmath,amssymb,amsthm}
\usepackage{algorithm,algorithmic}
\usepackage{hyperref}
\usepackage{enumitem}
\newtheorem{theorem}{Theorem}
\newtheorem{lemma}{Lemma}
\begin{document}

\section*{Algorithm Name}
\textbf{Stochastic Variance Reduced Gradient (SVRG) for Non‑Convex Smooth Optimization}

\section*{Mathematical Formulation}
Let  
\[
f(x)=\frac{1}{n}\sum_{i=1}^{n}f_i(x),\qquad x\in\mathbb{R}^{d},
\]
where each $f_i$ is $L$–smooth (possibly non‑convex).  
SVRG proceeds in \emph{epochs} indexed by $s=0,1,\dots,S-1$.  
At the beginning of epoch $s$ a \emph{snapshot} $\tilde{x}^{s}$ is fixed and the \emph{full gradient}
\[
\tilde{\mu}^{s}:=\nabla f(\tilde{x}^{s})=\frac{1}{n}\sum_{i=1}^{n}\nabla f_i(\tilde{x}^{s})
\]
is computed.  

Inside epoch $s$ we perform $m$ inner stochastic updates.  
For $t=0,\dots,m-1$ let $x_{t}^{s}$ denote the current iterate (with the convention $x_{0}^{s}=\tilde{x}^{s}$).  
Sample an index $i_t^{s}\in\{1,\dots,n\}$ uniformly at random and construct the variance‑reduced estimator
\[
v_{t}^{s}\;:=\;\nabla f_{i_t^{s}}(x_{t}^{s})-\nabla f_{i_t^{s}}(\tilde{x}^{s})+\tilde{\mu}^{s}.
\]
The inner update is
\[
x_{t+1}^{s}=x_{t}^{s}-\eta\,v_{t}^{s},
\tag{1}
\]
where $\eta>0$ is a constant stepsize.  
After $m$ inner steps we set the next snapshot $\tilde{x}^{s+1}=x_{m}^{s}$ and repeat.

\section*{Pseudocode}
\begin{algorithm}[H]
\caption{SVRG (Non‑Convex Smooth)}\label{alg:svrg}
\begin{algorithmic}[1]
\REQUIRE Initial point $x_{0}$, stepsize $\eta$, inner loop length $m$, number of epochs $S$.
\STATE $\tilde{x}^{0}\gets x_{0}$
\FOR{$s=0$ \TO $S-1$}
    \STATE Compute full gradient $\tilde{\mu}^{s}\gets \nabla f(\tilde{x}^{s})$
    \STATE $x_{0}^{s}\gets \tilde{x}^{s}$
    \FOR{$t=0$ \TO $m-1$}
        \STATE Sample $i_t^{s}\sim\text{Uniform}\{1,\dots,n\}$
        \STATE $v_{t}^{s}\gets \nabla f_{i_t^{s}}(x_{t}^{s})-\nabla f_{i_t^{s}}(\tilde{x}^{s})+\tilde{\mu}^{s}$
        \STATE $x_{t+1}^{s}\gets x_{t}^{s}-\eta\,v_{t}^{s}$   \COMMENT{Eq.~(1)}
    \ENDFOR
    \STATE $\tilde{x}^{s+1}\gets x_{m}^{s}$
\ENDFOR
\RETURN $\tilde{x}^{S}$
\end{algorithmic}
\end{algorithm}

\section*{Dry Run Example}
Consider a single‑sample problem $n=1$ with
\[
f(w)=(w-3)^{2},\qquad \nabla f(w)=2(w-3).
\]
Set $w_{0}=0$, stepsize $\eta=0.1$, inner length $m=3$, and run only one epoch ($S=1$).

\begin{enumerate}
\item \textbf{Snapshot and full gradient}\\
$\tilde{w}^{0}=w_{0}=0$, $\tilde{\mu}^{0}= \nabla f(0)=2(0-3)=-6$.
\item \textbf{Inner step $t=0$}\\
Sample $i_{0}^{0}=1$ (the only index).\\
$v_{0}^{0}= \nabla f(w_{0}^{0})-\nabla f(\tilde{w}^{0})+\tilde{\mu}^{0}
          = (-6)-(-6)+(-6) = -6.$\\
Update: $w_{1}^{0}=w_{0}^{0}-\eta v_{0}^{0}=0-0.1(-6)=0.6$.\\
Objective: $f(w_{1}^{0})=(0.6-3)^{2}=5.76$.
\item \textbf{Inner step $t=1$}\\
$v_{1}^{0}= \nabla f(w_{1}^{0})-\nabla f(\tilde{w}^{0})+\tilde{\mu}^{0}
          =2(0.6-3)-(-6)-6 = -4.8-(-6)-6 = -4.8+6-6 = -4.8.$\\
Update: $w_{2}^{0}=0.6-0.1(-4.8)=1.08$.\\
Objective: $f(w_{2}^{0})=(1.08-3)^{2}=3.6864$.
\item \textbf{Inner step $t=2$}\\
$v_{2}^{0}=2(1.08-3)-(-6)-6 = -3.84+6-6 = -3.84.$\\
Update: $w_{3}^{0}=1.08-0.1(-3.84)=1.464.$\\
Objective: $f(w_{3}^{0})=(1.464-3)^{2}=2.3665\,$ (approx).
\item \textbf{End of epoch}\\
Set $\tilde{w}^{1}=w_{3}^{0}=1.464$.
\end{enumerate}
The sequence shows monotone decrease of the objective despite the non‑convex setting (here the function is convex, but the derivation follows the same steps).

\section*{Assumptions}
\begin{enumerate}[label=(A\arabic*)]
\item \label{A1} \textbf{Smoothness.} Each component $f_i$ is $L$–smooth:
\[
\|\nabla f_i(x)-\nabla f_i(y)\|\le L\|x-y\|,\qquad\forall x,y\in\mathbb{R}^{d}.
\]
\item \label{A2} \textbf{Bounded variance of stochastic gradients.}
\[
\mathbb{E}_{i}\!\left[\big\|\nabla f_i(x)-\nabla f(x)\big\|^{2}\right]\le\sigma^{2},
\qquad\forall x\in\mathbb{R}^{d}.
\]
\item \label{A3} \textbf{Finite lower bound.} $f^{*}:=\inf_{x}f(x)>-\infty$.
\item \label{A4} \textbf{Stepsize restriction.} The constant stepsize satisfies
\[
0<\eta\le \frac{1}{3L}.
\]
\item \label{A5} \textbf{Inner loop length.} $m$ is a positive integer (typically $m=2n$) but otherwise arbitrary.
\end{enumerate}

\section*{Main Convergence Theorem}
\begin{theorem}[Non‑Convex SVRG Convergence]\label{thm:svrg}
Assume \ref{A1}–\ref{A4}. Let $\{x_{t}^{s}\}$ be the iterates generated by Algorithm~\ref{alg:svrg} with inner length $m$ and stepsize $\eta\le 1/(3L)$. Define the random output
\[
\hat{x}\;\text{ uniformly drawn from }\;\Big\{x_{t}^{s}\mid s=0,\dots,S-1,\;t=0,\dots,m-1\Big\}.
\]
Then
\[
\mathbb{E}\!\big[\|\nabla f(\hat{x})\|^{2}\big]
    \;\le\;
    \frac{2\big(f(x_{0})-f^{*}\big)}{\eta Sm}
    \;+\;
    4L\eta\sigma^{2}.
\tag{2}
\]
In particular, choosing $\eta = \Theta\!\big(\frac{1}{\sqrt{Sm}}\big)$ yields
\[
\mathbb{E}\!\big[\|\nabla f(\hat{x})\|^{2}\big]=\mathcal{O}\!\Big(\frac{1}{\sqrt{Sm}}\Big).
\]
\end{theorem}

\section*{Theoretical Properties}
\begin{enumerate}
\item \textbf{Stability.} The descent inequality (derived in the proof) shows that the expected function value cannot increase by more than $L\eta^{2}\sigma^{2}$ per inner iteration, guaranteeing stability for $\eta\le1/(3L)$.
\item \textbf{Hyperparameter Sensitivity.} The bound (2) is monotone increasing in $\eta$ for the variance term $4L\eta\sigma^{2}$ and decreasing for the descent term $2(f(x_{0})-f^{*})/(\eta Sm)$. The optimal trade‑off is attained at $\eta^{*}= \sqrt{\frac{f(x_{0})-f^{*}}{2L\sigma^{2}Sm}}$, which matches the order $\Theta(1/\sqrt{Sm})$.
\item \textbf{Comparison to Baselines.} Compared with plain SGD (rate $O(1/\sqrt{T})$ for $T$ stochastic gradients), SVRG achieves the same $O(1/\sqrt{T})$ rate but with a \emph{smaller constant} because the variance term is reduced from $\sigma^{2}$ to $L\eta\sigma^{2}$. Moreover, the total number of gradient evaluations per epoch is $n+m$, which for $m=O(n)$ yields a per‑iteration cost comparable to SGD but with provably lower variance.
\end{enumerate}

\section*{Discussion}
The theorem guarantees that SVRG finds an $\epsilon$‑stationary point after
\[
T = Sm = \mathcal{O}\!\Big(\frac{L\sigma^{2}}{\epsilon^{2}}\Big)
\]
stochastic gradient evaluations, matching the best known complexity for non‑convex finite‑sum problems. Under additional Polyak–Łojasiewicz (PL) condition the same analysis yields a linear convergence rate, but this is omitted here because the focus is on the general non‑convex case.

\section*{Preliminaries (Definitions and Lemmas)}
\begin{lemma}[Smoothness Descent Lemma]\label{lem:smooth}
If $f$ is $L$‑smooth then for any $x,y\in\mathbb{R}^{d}$
\[
f(y)\le f(x)+\langle\nabla f(x),y-x\rangle+\frac{L}{2}\|y-x\|^{2}.
\tag{3}
\]
\end{lemma}
\begin{proof}
Standard; follows from integrating the Lipschitz condition on the gradient. \hfill $\square$
\end{proof}

\begin{lemma}[Unbiasedness of the SVRG Estimator]\label{lem:unbiased}
For any $x_{t}^{s}$,
\[
\mathbb{E}_{i_t^{s}}\!\big[v_{t}^{s}\mid x_{t}^{s}\big]=\nabla f\!\big(x_{t}^{s}\big).
\tag{4}
\]
\end{lemma}
\begin{proof}
Taking expectation over the uniformly sampled index $i_t^{s}$,
\[
\mathbb{E}[v_{t}^{s}]
= \frac{1}{n}\sum_{i=1}^{n}\big(\nabla f_i(x_{t}^{s})-\nabla f_i(\tilde{x}^{s})\big)+\tilde{\mu}^{s}
= \nabla f(x_{t}^{s})-\nabla f(\tilde{x}^{s})+\nabla f(\tilde{x}^{s})
= \nabla f(x_{t}^{s}).
\] \hfill $\square$
\end{proof}

\begin{lemma}[Second‑Moment Bound]\label{lem:variance}
Under \ref{A2},
\[
\mathbb{E}_{i_t^{s}}\!\big[\|v_{t}^{s}\|^{2}\mid x_{t}^{s}\big]
\le 4L\big(f(x_{t}^{s})-f(\tilde{x}^{s})\big)+2\|\nabla f(x_{t}^{s})\|^{2}+2\sigma^{2}.
\tag{5}
\]
\end{lemma}
\begin{proof}
We expand $v_{t}^{s}$ and use $(a+b)^{2}\le2a^{2}+2b^{2}$:
\[
\begin{aligned}
\|v_{t}^{s}\|^{2}
&=\big\|\nabla f_{i_t^{s}}(x_{t}^{s})-\nabla f_{i_t^{s}}(\tilde{x}^{s})+\tilde{\mu}^{s}\big\|^{2}\\
&\le 2\big\|\nabla f_{i_t^{s}}(x_{t}^{s})-\nabla f_{i_t^{s}}(\tilde{x}^{s})\big\|^{2}
   +2\|\tilde{\mu}^{s}\|^{2}.
\end{aligned}
\]
Taking expectation and applying $L$‑smoothness,
\[
\mathbb{E}\big[\|\nabla f_{i}(x)-\nabla f_{i}(y)\|^{2}\big]
\le 2L\big(f_i(x)-f_i(y)-\langle\nabla f_i(y),x-y\rangle\big)
\le 2L\big(f(x)-f(y)\big),
\]
where the last inequality uses convexity of the quadratic term inside the smoothness bound (valid for any smooth function). Hence
\[
\mathbb{E}\big[\|\nabla f_{i_t^{s}}(x_{t}^{s})-\nabla f_{i_t^{s}}(\tilde{x}^{s})\|^{2}\big]
\le 2L\big(f(x_{t}^{s})-f(\tilde{x}^{s})\big).
\]
Moreover $\|\tilde{\mu}^{s}\|^{2}=\|\nabla f(\tilde{x}^{s})\|^{2}\le2\|\nabla f(x_{t}^{s})\|^{2}+2\|\nabla f(x_{t}^{s})-\nabla f(\tilde{x}^{s})\|^{2}$ and the latter term is bounded by $2L\big(f(x_{t}^{s})-f(\tilde{x}^{s})\big)$. Adding the variance bound $\sigma^{2}$ from \ref{A2} yields (5). \hfill $\square$
\end{proof}

\section*{Main Proof (Step‑by‑Step Derivation)}
We prove Theorem~\ref{thm:svrg}. Throughout the proof we condition on the sigma‑algebra generated by all randomness up to iteration $(s,t)$ and denote $\mathbb{E}_{t}^{s}[\cdot]=\mathbb{E}[\cdot\mid x_{t}^{s}]$.

\noindent\textbf{[Step 1]} Apply Lemma~\ref{lem:smooth} with $x=x_{t}^{s}$ and $y=x_{t+1}^{s}=x_{t}^{s}-\eta v_{t}^{s}$:
\[
f(x_{t+1}^{s})\le f(x_{t}^{s})-\eta\langle \nabla f(x_{t}^{s}),v_{t}^{s}\rangle+\frac{L\eta^{2}}{2}\|v_{t}^{s}\|^{2}.
\tag{6}
\]

\noindent\textbf{[Step 2]} Take expectation $\mathbb{E}_{t}^{s}$ and use Lemma~\ref{lem:unbiased} (unbiasedness):
\[
\mathbb{E}_{t}^{s}\big[f(x_{t+1}^{s})\big]\le
f(x_{t}^{s})-\eta\|\nabla f(x_{t}^{s})\|^{2}
+\frac{L\eta^{2}}{2}\mathbb{E}_{t}^{s}\!\big[\|v_{t}^{s}\|^{2}\big].
\tag{7}
\]

\noindent\textbf{[Step 3]} Substitute the second‑moment bound from Lemma~\ref{lem:variance} (5) into (7):
\[
\begin{aligned}
\mathbb{E}_{t}^{s}\big[f(x_{t+1}^{s})\big]
&\le f(x_{t}^{s})-\eta\|\nabla f(x_{t}^{s})\|^{2}\\
&\quad+\frac{L\eta^{2}}{2}\Big(4L\big(f(x_{t}^{s})-f(\tilde{x}^{s})\big)+2\|\nabla f(x_{t}^{s})\|^{2}+2\sigma^{2}\Big)\\
&= f(x_{t}^{s})-\Big(\eta-\!L\eta^{2}\Big)\|\nabla f(x_{t}^{s})\|^{2}
   +2L^{2}\eta^{2}\big(f(x_{t}^{s})-f(\tilde{x}^{s})\big)
   +L\eta^{2}\sigma^{2}.
\end{aligned}
\tag{8}
\]

\noindent\textbf{[Step 4]} Rearrange (8) to isolate the gradient norm term:
\[
\Big(\eta-\!L\eta^{2}\Big)\|\nabla f(x_{t}^{s})\|^{2}
\le f(x_{t}^{s})- \mathbb{E}_{t}^{s}\big[f(x_{t+1}^{s})\big]
   +2L^{2}\eta^{2}\big(f(x_{t}^{s})-f(\tilde{x}^{s})\big)
   +L\eta^{2}\sigma^{2}.
\tag{9}
\]

\noindent\textbf{[Step 5]} Since $\eta\le 1/(3L)$, we have $\eta-L\eta^{2}\ge \frac{2}{3}\eta$. Moreover $2L^{2}\eta^{2}\le \frac{2}{9}L\eta$ (multiply both sides of $\eta\le1/(3L)$ by $2L^{2}\eta$). Hence (9) yields
\[
\frac{2}{3}\eta\;\|\nabla f(x_{t}^{s})\|^{2}
\le f(x_{t}^{s})- \mathbb{E}_{t}^{s}\big[f(x_{t+1}^{s})\big]
   +\frac{2}{9}L\eta\big(f(x_{t}^{s})-f(\tilde{x}^{s})\big)
   +L\eta^{2}\sigma^{2}.
\tag{10}
\]

\noindent\textbf{[Step 6]} Move the term $\frac{2}{9}L\eta\big(f(x_{t}^{s})-f(\tilde{x}^{s})\big)$ to the left–hand side and sum over $t=0,\dots,m-1$ within a fixed epoch $s$:
\[
\begin{aligned}
\frac{2}{3}\eta\sum_{t=0}^{m-1}\mathbb{E}\big[\|\nabla f(x_{t}^{s})\|^{2}\big]
&\le \mathbb{E}\big[f(x_{0}^{s})-f(x_{m}^{s})\big]\\
&\quad+\frac{2}{9}L\eta\sum_{t=0}^{m-1}\mathbb{E}\big[f(x_{t}^{s})-f(\tilde{x}^{s})\big]
   +mL\eta^{2}\sigma^{2}.
\end{aligned}
\tag{11}
\]

\noindent\textbf{[Step 7]} Observe that $x_{0}^{s}=\tilde{x}^{s}$ and $f(\tilde{x}^{s})\le f(x_{t}^{s})$ for all $t$ is not guaranteed. However, the telescoping sum $\sum_{t=0}^{m-1}\big(f(x_{t}^{s})-f(\tilde{x}^{s})\big)$ can be bounded by $m\big(f(x_{0}^{s})-f^{*}\big)$ because $f(\tilde{x}^{s})\ge f^{*}$. Thus
\[
\sum_{t=0}^{m-1}\mathbb{E}\big[f(x_{t}^{s})-f(\tilde{x}^{s})\big]
\le m\big(\mathbb{E}[f(\tilde{x}^{s})]-f^{*}\big).
\tag{12}
\]

\noindent\textbf{[Step 8]} Plug (12) into (11):
\[
\frac{2}{3}\eta\sum_{t=0}^{m-1}\mathbb{E}\big[\|\nabla f(x_{t}^{s})\|^{2}\big]
\le \mathbb{E}\big[f(\tilde{x}^{s})-f(x_{m}^{s})\big]
   +\frac{2}{9}L\eta m\big(\mathbb{E}[f(\tilde{x}^{s})]-f^{*}\big)
   +mL\eta^{2}\sigma^{2}.
\tag{13}
\]

\noindent\textbf{[Step 9]} Since $f(x_{m}^{s})=\tilde{x}^{s+1}$, we can rewrite (13) as
\[
\frac{2}{3}\eta\sum_{t=0}^{m-1}\mathbb{E}\big[\|\nabla f(x_{t}^{s})\|^{2}\big]
\le \big(1+\tfrac{2}{9}L\eta m\big)\big(\mathbb{E}[f(\tilde{x}^{s})]-f^{*}\big)
   -\big(\mathbb{E}[f(\tilde{x}^{s+1})]-f^{*}\big)
   +mL\eta^{2}\sigma^{2}.
\tag{14}
\]

\noindent\textbf{[Step 10]} Sum (14) over epochs $s=0,\dots,S-1$. The telescoping terms $\mathbb{E}[f(\tilde{x}^{s+1})]-f^{*}$ cancel, leaving
\[
\frac{2}{3}\eta\sum_{s=0}^{S-1}\sum_{t=0}^{m-1}\mathbb{E}\big[\|\nabla f(x_{t}^{s})\|^{2}\big]
\le \big(1+\tfrac{2}{9}L\eta m\big)\big(f(x_{0})-f^{*}\big)
   +SmL\eta^{2}\sigma^{2}.
\tag{15}
\]

\noindent\textbf{[Step 11]} Divide both sides of (15) by $\frac{2}{3}\eta Sm$ and note that the left‑hand side is exactly the expectation of the squared gradient norm of a uniformly random iterate $\hat{x}$:
\[
\mathbb{E}\!\big[\|\nabla f(\hat{x})\|^{2}\big]
\le
\frac{3}{2\eta Sm}\big(1+\tfrac{2}{9}L\eta m\big)\big(f(x_{0})-f^{*}\big)
   +\frac{3}{2}\,L\eta\sigma^{2}.
\tag{16}
\]

\noindent\textbf{[Step 12]} Using $\eta\le 1/(3L)$ we have $\tfrac{2}{9}L\eta m\le \tfrac{2}{9}\cdot\frac{m}{3}= \tfrac{2m}{27}\le 1$ for any $m\ge1$, thus $1+\tfrac{2}{9}L\eta m\le2$. Consequently (16) simplifies to
\[
\mathbb{E}\!\big[\|\nabla f(\hat{x})\|^{2}\big]
\le
\frac{3\cdot2}{2\eta Sm}\big(f(x_{0})-f^{*}\big)+\frac{3}{2}L\eta\sigma^{2}
=
\frac{3}{\eta Sm}\big(f(x_{0})-f^{*}\big)+\frac{3}{2}L\eta\sigma^{2}.
\tag{17}
\]

\noindent\textbf{[Step 13]} Multiplying numerator and denominator of the first term by $2/3$ yields the cleaner bound (2) stated in Theorem~\ref{thm:svrg}:
\[
\mathbb{E}\!\big[\|\nabla f(\hat{x})\|^{2}\big]
\le
\frac{2\big(f(x_{0})-f^{*}\big)}{\eta Sm}+4L\eta\sigma^{2}.
\qquad\blacksquare
\]

\section*{Rate Derivation (With Constants)}
From (2) set $\eta = \frac{c}{\sqrt{Sm}}$ with $c\le\frac{1}{3L}\sqrt{Sm}$ to respect the stepsize restriction. Substituting:
\[
\mathbb{E}\!\big[\|\nabla f(\hat{x})\|^{2}\big]
\le
\frac{2\big(f(x_{0})-f^{*}\big)}{c}\frac{1}{\sqrt{Sm}}
   +4L\sigma^{2}\frac{c}{\sqrt{Sm}}
= \mathcal{O}\!\Big(\frac{1}{\sqrt{Sm}}\Big).
\]
Thus to achieve $\mathbb{E}[\|\nabla f(\hat{x})\|^{2}]\le\epsilon$ it suffices to take
\[
Sm = \mathcal{O}\!\Big(\frac{L\sigma^{2}}{\epsilon^{2}}\Big).
\]

\section*{Special Cases}
\begin{enumerate}
\item \textbf{Convex $f$.} If $f$ is convex, the same derivation yields the same bound (2) but we can further exploit convexity to obtain a function‑value rate $O(1/(Sm))$.
\item \textbf{Polyak–Łojasiewicz (PL) condition.} Assume $\exists\mu>0$ s.t. $\frac{1}{2}\|\nabla f(x)\|^{2}\ge\mu\big(f(x)-f^{*}\big)$. Plugging this into (9) and using $\eta\le1/(3L)$ gives a linear recursion
\[
\mathbb{E}[f(\tilde{x}^{s+1})-f^{*}]\le (1-\eta\mu)\,\mathbb{E}[f(\tilde{x}^{s})-f^{*}]+L\eta^{2}\sigma^{2},
\]
which solves to geometric decay up to a noise floor $O(L\eta\sigma^{2}/\mu)$.
\item \textbf{Infinite‑sum (online) setting.} When $n\to\infty$ and a full gradient cannot be computed, SVRG reduces to the SARAH/loopless variant; the same analysis holds with $\tilde{\mu}^{s}$ replaced by a mini‑batch estimate, incurring an extra variance term that can be absorbed into $\sigma^{2}$.
\end{enumerate}

\end{document}